\begin{table}[H]
\centering\small
\caption{Insurance quality of the major hubs' demand. Import demand, what pulls a client's exports, is steadiest and growing for the United States; China's grows but is twice as volatile. Ranking hubs by GDP growth instead (last column) misleads, placing China first, because China grew as a producer rather than a consumer.}
\label{tab:partnerquality}
\begin{tabularx}{\textwidth}{@{}lYYYY@{}}
\toprule
Hub & Import-demand & Volatility & Quality & GDP-growth \\
 & growth, \% & of demand, \% & (mean/vol.) & ratio \\ \midrule
United States & 4.2 & 6.7 & 0.62 & 2.57 \\
European Union & 3.1 & 10.2 & 0.30 & 1.25 \\
China & 5.6 & 13.0 & 0.43 & 3.77 \\
Japan & 1.2 & 12.2 & 0.10 & 1.01 \\
\bottomrule
\end{tabularx}
\\[2pt]\footnotesize{\emph{Source:} IMF IMTS (import demand, total goods imports from the world) and WEO (real GDP growth), 2010--2024 excluding the 2020 pandemic outlier. Import demand for the European Union aggregates its ten largest members. Quality is mean annual import growth divided by its standard deviation (computed from unrounded series); the GDP-growth ratio is mean annual real GDP growth divided by its standard deviation over the same window.}
\end{table}
